\section{Motivation}
\label{sec:intro-motivation}

During the three years in which this thesis was written, the practical meaning
of \emph{machine-assisted reasoning} changed. In my own work, general-purpose
artificial intelligence (AI) assistants moved from brief, isolated exchanges
to sustained sessions involving sizable programs and formal developments,
proposing multi-step solutions and responding repeatedly to feedback from
compilers and proof assistants. These tools, initially peripheral to the
research process, became part of its daily practice. This transition gives the
subject of this manuscript---trustworthy reasoning---an urgency that was less
visible when the work began.

Computers did not begin assisting reasoning with large language models (LLMs).
They have long performed calculations, explored large spaces of possibilities,
and checked arguments that would be exhausting to follow by hand. What is new
is the breadth and apparent fluency of what a general-purpose system can
propose. LLMs now produce source code, technical explanations, proof scripts,
reviews, hypotheses, and research plans in forms that resemble the products of
human reasoning. They participate through the same languages in which
researchers formulate ideas and communicate results.

The underlying systems have progressed rapidly enough that any detailed
forecast risks becoming obsolete before it is printed. Nevertheless, the
direction is already visible.
Systems are being designed not only to answer questions but to participate in
hypothesis generation, experimental planning, programming, and other parts of
the scientific workflow~\cite{stanfordAIIndex2026,gottweis2026coscientist}.
Whether or not the strongest predictions about autonomous research are
realized, it is prudent to expect research to accommodate a growing volume of
machine-generated claims and artifacts.

This expansion of capability does not remove the problem of reliability; it
makes it more pressing. Striking results in mathematics, coding, and scientific
question answering coexist with brittle behavior on apparently simpler tasks
and with poor performance on end-to-end scientific
replication~\cite{stanfordAIIndex2026}. A generated answer may be correct,
subtly wrong, or correct for reasons unrelated to the explanation that
accompanies it. Fluency removes the surface cues by which an unreliable answer
used to announce itself, so the three cases now arrive equally well presented.
As the cost of producing plausible answers falls, justifications that deserve
confidence gain in importance.

Leonardo de Moura recently captured this reciprocal promise in three
sentences~\cite{deMoura2026floc}:
\begin{quotation}[Leonardo de Moura][2026][The Lean Theorem Prover: Design, Evolution, and Impact]
  We have never been more relevant.

  Formal verification makes AI trustworthy.

  AI makes formal verification cost-effective.
\end{quotation}

These lines capture the perspective adopted in this thesis: formal verification
can provide an independent criterion for accepting AI-generated artifacts,
while AI can reduce the effort required to construct formally checked ones.
This reciprocity renews the relevance of formal methods. At their broadest,
formal methods are not synonymous with the familiar slogan of ``proving
software correct.'' They provide languages and calculi in which artifacts,
assumptions, and arguments can be made precise. The artifact may be a program,
a hardware design, a communication protocol, a system model, or a mathematical
theory; the claim may concern safety, consistency, refinement, equivalence,
security, or some other property. What unites these applications is the attempt
to replace an appeal to plausibility with an explicit argument whose individual
steps can be checked.

Formalization does not create certainty from nothing. A theorem can faithfully
establish the wrong requirement, or the right one under hypotheses that nothing
in practice satisfies; an axiom can be inappropriate; a model can omit
an important aspect of the world; and the software that performs the checking
is itself part of the trust placed in the result. Formal methods are valuable
because they make these dependencies more visible and more sharply delimited.
They turn the question ``Do we trust this answer?'' into more concrete
questions: what exactly was stated, from which assumptions was it derived, what
mechanism checked the derivation, and which components remain in the trusted
base?

Generative AI and formal methods can therefore be complementary. A human, an
LLM, or a search procedure may suggest an invariant, a program, or a proof. A
formal system can then reject the suggestion or certify that it satisfies a
precisely stated criterion. The creative component may be probabilistic and
opaque while the acceptance criterion remains deterministic and inspectable.
Recent systems for mathematical reasoning illustrate this architecture:
AlphaProof, for example, couples learned proof search with the Lean proof
assistant, so that successful searches yield formal proofs checked by Lean's
kernel~\cite{hubert2026alphaproof}. The result is trustworthy because acceptance
does not depend on trusting the model's account of its own reasoning.

This thesis adopts that separation as a guiding principle. It contributes no new
automation; it works on the links that make automation acceptable, so that
increasingly powerful proof-producing tools---including AI-assisted ones---can be
used without making the validity of a result depend on the opacity of the tool
that proposed it. \Cref{sec:intro-llm-use} separately accounts for how LLMs were
used in producing this work.

\subsection{Two modes of machine-assisted reasoning}
\label{subsec:intro-itp-atp}

Computer-assisted proof is commonly organized around two complementary types of
systems. An \emph{automated theorem prover} (ATP) receives a conjecture and
searches for a proof with little or no guidance from its user. Modern
satisfiability modulo theories (SMT) solvers, such as
Z3~\cite{demoura2008z3} and cvc5~\cite{barbosa2022cvc5}, are particularly
effective ATPs for verification conditions: they combine propositional search
with specialized procedures for equality, arithmetic, arrays, and other useful
theories. Their main attraction is throughput: once obligations have been
translated into the solver's language, large batches can be processed
automatically.

An \emph{interactive theorem prover} (ITP), or proof assistant, instead provides
an environment in which definitions, statements, and proofs are developed
together. The user may guide the argument step by step or invoke substantial
automation, but every successful proof is ultimately represented by a proof term
checked by a comparatively small kernel. Lean~4~\cite{deMouraUllrich2021Lean4},
used throughout this thesis, combines this kernel architecture with dependent
types, an extensible tactic and metaprogramming framework, and the large
mathematical library Mathlib~\cite{mathlib4}. Its main attractions are
expressiveness, direct access to the proof state, and the possibility of
independently rechecking the resulting theorem. The corresponding cost is that
proofs that lie beyond the available automation require insight and
maintenance.

The distinction is therefore not simply ``automatic versus manual.'' ITPs
routinely run powerful tactics, and ATPs may emit proof certificates that
another program can check. Rather, the difference lies in how proof search,
interaction, and trust are organized.

\subsection{The bridge is part of the proof}
\label{subsec:intro-bridge-trust}

Using either kind of prover as a backend for an existing formalism requires a
translation. The source and target languages---in this thesis, B and SMT-LIB---may differ
in their type systems, primitive constructions, mathematical foundations, and
representation of binders. Connecting them is therefore not merely a matter of
printing the same formula in another syntax. The translation must choose
representations, recover the information required by the target, and preserve
the meaning of the source statement.

This bridge belongs to the assurance argument. A flawless prover cannot repair
a translation that changes the meaning of its input: an invalid source
statement may become a valid target statement, causing the backend to solve the
wrong problem correctly. Likewise, checking a target proof establishes only the
target statement unless its relationship with the source statement is itself
justified. The more capable the backend becomes, the more important it is not to
hide this boundary behind the success message of a tool.

This thesis originates in the tension between \emph{automation} and
\emph{assurance}.
Addressing it requires more than a new collection of prover heuristics: it
requires a mathematical account of the languages being connected, an
implementation architecture in which semantic claims can be stated and checked,
and a clear separation between what has been proved and what remains trusted.
The contributions presented next address these requirements while pursuing the
automated and interactive routes in parallel.
